% Conclusion
% Target: 0.5 pages

We presented a systematic comparison of speech science acoustic features versus deep learning embeddings for emotion-augmented knowledge distillation in small language models.

\textbf{Key Results:}
\begin{itemize}
    \item 46D acoustic features achieve \textbf{40.45\% validation loss reduction}
    \item Outperform 256D WavLM embeddings (27.82\%) by \textbf{17.5\%}
    \item 8.06$\times$ dimensional efficiency with superior performance
    \item Validated across 100, 500, and 1000 sample scales
\end{itemize}

\textbf{Contributions:}
\begin{enumerate}
    \item First systematic study of speech science vs. deep learning for emotion-augmented LLMs
    \item Design of interpretable 46D acoustic feature set
    \item Empirical demonstration that domain knowledge outperforms learned representations
    \item Analysis of feature importance and generalization
\end{enumerate}

\textbf{Impact:}
This work challenges the conventional wisdom that deep learning always outperforms hand-crafted features, demonstrating the value of speech science principles for practical deployment.

\textbf{Future Work:}
We plan to extend this approach to real human speech, larger scales, and multi-lingual settings, establishing acoustic features as a production-ready solution for emotion-aware LLMs.
